\section{Methods}
\subsection{Study population and clinical assessment}
We analysed longitudinal records from a Chinese observational treat-to-target project at CAMS and Huashan. The supplied protocol specified Hanifin--Rajka diagnosis and eligibility for systemic treatment. Actual screening implementation and recruitment dates were not recoverable from a verified screening register. Systemic treatment was clinician directed. Initial regimens were dupilumab, abrocitinib 100 mg or abrocitinib 200 mg; later dose, interval and regimen changes were recorded separately. Clinical comparisons included patients aged at least 12 years with an unambiguous initial regimen. Adult Dermatology Life Quality Index (DLQI) analyses required age at least 18.

Clinical outcomes were the recorded EASI, peak pruritus numerical rating scale (PP-NRS), Investigator Global Assessment (IGA), DLQI, BSA and SCORAD, with IgE and eosinophil measurements. Reporting visits were W0/W2/W4/W12/W24 for PP-NRS, W0/W12/W24 for IgE and W0/W4/W12/W24 for other measures. Records were linked by patient and nominal visit. The exact IGA and questionnaire versions, PP-NRS recall window and calendar-day clinical/serum matching require source confirmation. Source discrepancies were audited without substituting imputed clinical outcomes.

\subsection{Clinical framework, targets and treatment pathways}
Project materials describe evidence review, a two-round questionnaire procedure and a planned agreement criterion of at least 75\%. They report 49 returned questionnaires, without original ballots or round-specific agreement results. We transcribed the September framework and operationalized its targets on recorded scores (Fig.~\ref{fig:Figure1}). Its exact approved and prospectively implemented version remains unverified; this analysis does not validate a new consensus or the effectiveness of the management strategy.

The week-2 target was PP-NRS $\leq4$. At week 4, the physician target was EASI50 or an IGA decrease of at least two points, and the patient target was PP-NRS $\leq4$ or adult DLQI $\leq10$. At week 12, these were EASI75 or IGA 0/1, and PP-NRS 0/1 or adult DLQI $\leq4$. Week 24 additionally allowed EASI $\leq5$ in the physician domain. Both domains were required for the combined target. These investigator-defined targets were kept distinct from newer low-disease-activity and remission definitions.\citep{merola2026iec}

Each rate used its own evaluable denominator and Wilson 95\% confidence interval. EASI percentage responses required positive observed baseline EASI. A PP-NRS decrease of at least four points required observed baseline PP-NRS $\geq4$ and an observed follow-up score; lower or missing baseline values were excluded from that endpoint. Missing components remained unknown unless another observed component determined the logical result. Below age 18, adult DLQI was treated as unknown: a positive itch component establishes the patient target, whereas a negative itch component alone leaves it unknown. Group comparisons retained initial treatment after subsequent changes. Separate maintenance analyses classified week-12 response and later reduction pathways, then examined observed week-24 EASI75 with explicit missing-outcome bounds. Reduction was not reassigned to an earlier visit; uncertain assessment/decision timing was retained. These descriptive pathways address treatment changes without treating a response-selected maintenance group as randomized.\citep{iche9r1}

\subsection{Serum measurements}
Source laboratory reports described Luminex and MSD serum assays, including the LXSAHM-13 panel on Luminex-200 and K15067M-1/K15067M-2 kits on SQ120, using 50-$\mu$L serum and twofold dilution. All 23 measured markers were included. Analyses used positive numeric integration-export concentrations recorded in pg/mL, transformed to log$_2$ once without an additional dilution correction. The twofold laboratory dilution describes sample preparation; its correction in the export and validated quantification limits require source confirmation. Censoring strings and nonpositive values were not treated as measured positive concentrations. An IL-5 calibration-bound sensitivity analysis and detailed assay provenance are provided in the Supplementary Methods; software calibration limits were distinguished from validated quantification limits.

Molecular changes were calculated within patient for W0--W12, W12--W24 and W0--W24 before summarization. Three-visit displays additionally held the patient set fixed within each marker and treatment. This separates individual change from shifts in the patients observed at each visit.

\subsection{EASI estimation and validation}
Eligible model visits required age at least 12, observed EASI and BSA, and positive observed CCL17, CCL26 and IL-5. This three-marker panel was fixed from earlier exploratory work. Development used W0/W12; W24 was excluded from fitting and selection. As W24 had been examined in previous project analyses, it represents within-cohort temporal application, not untouched external validation.

A repeated five-fold patient-level outer validation with three-fold inner validation selected among three common structures and ridge penalties. All records from a patient stayed in one outer fold, preprocessing used training data only, and patients received equal total fitting weight. Full-development selection yielded a formula in which BSA multiplies a molecular component. Its rounded display is
\begin{align*}
\widehat{\mathrm{EASI}}\ \approx\ 
\operatorname{clip}_{[0,72]}\Bigl\{\frac{\mathrm{BSA}}{100}
\bigl[&51.27+0.98\log_2(\mathrm{CCL17})\\
&+0.96\log_2(\mathrm{CCL26})
-0.83\log_2(\mathrm{IL\mbox{-}5})\bigr]\Bigr\}.
\end{align*}
BSA is entered from 0 to 100 and concentrations use the development export's recorded pg/mL numerics, with the logarithm applied once. Here, clip bounds the result to 0--72, and log$_2$ is the base-two logarithm. The project label Skin inflammation score refers here to this hybrid EASI estimate, not a separately calibrated blood-only inflammation scale. The calculator retains full-precision parameters; the displayed approximation differed from final-equation predictions by less than 0.08 EASI points in this dataset. Threshold decisions use the precise calculator. Treatment, stage, time and previous scores are not inputs. Reported primary performance evaluates the nested selection procedure; it is not the apparent fit of this final equation.

Primary assessment used MAE with patient bootstrap intervals conditional on the saved held-out predictions; these intervals did not repeat model selection or fitting. This internal validation does not include the earlier selection of the marker panel. Additional measures were RMSE, signed error, Spearman correlation, Lin's concordance correlation and proportions within one or two EASI points. Repetition-specific metrics were calculated before averaging. Severity strata and centre/treatment holdouts assessed where performance varied. A separate refitting sensitivity removed development visits with complete regional scores whose sum differed from recorded total EASI by more than 0.15 points. It retained the original patient folds and outcomes and repeated inner selection; it tested exclusion robustness without correcting either clinical reference. Monitoring used estimated and observed changes within the same patient, repetition and fold; response classification used the corresponding estimated or observed scores at both ends of the interval. Matched controls included BSA alone and the three markers alone. Primary paired comparisons used the nested procedure; fixed-structure comparisons were labelled separately. Published p-EASI coefficients were applied unchanged as a source-table sensitivity analysis, using local IL2RA as an unbridged soluble IL-2 receptor proxy and retaining raw and matched-clipping results. Different inputs and unverified assay equivalence preclude a model-superiority claim.\citep{thijs2017,thijs2019,bakker2020} PP-NRS, adult DLQI and recorded IGA were fitted independently with target-appropriate regression or ordinal models using the same restricted inputs.

\subsection{Statistical analysis}
Patients were the independent units. Overall clinical trajectories showed means with 95\% intervals; treatment-specific observations showed medians and interquartile ranges. Continuous treatment comparisons used baseline-outcome adjustment, treatment-by-visit terms, age, sex, BMI and centre, generalized overlap weights and patient-cluster robust covariance. Weights target overlap in baseline treatment probabilities. Reported outcome contrasts are baseline-weighted associations among patients with that endpoint observed. Selective endpoint observation can therefore change the represented population and its covariate balance.\citep{lili2019} Weighted target contrasts used patient bootstrap intervals with propensity refitting. A W2 sensitivity additionally included baseline PP-NRS, comparing models on the same baseline-itch-complete cohort. Endpoint-specific balance, effective sample sizes and both weighted and unweighted missing-outcome bounds were retained; weighting did not correct informative follow-up loss. IgE/eosinophil models used log$(1+x)$ and back-transformation. The analyses estimate adjusted observational associations.

Unadjusted target comparisons used conditional Monte Carlo three-group tests and pairwise Fisher tests; the three pairwise contrasts within each endpoint/visit formed a Holm family. Molecular paired sign-flip tests used Benjamini--Hochberg adjustment over 23 markers, three treatments and three windows. Between-treatment molecular changes additionally adjusted for starting concentration and clinical baseline variables. Estimates based on fewer than five pairs were descriptive. Balance, effective sample size and common support were examined. Exact tests, numerical settings, model diagnostics, regional analyses and additional targets are detailed in the Supplementary Methods. Analyses used Python 3.9.6; dependency versions, code, data-access conditions and reproduction logs accompany the package. Reporting follows relevant STROBE and TRIPOD+AI domains.\citep{vonelm2007strobe,collins2024tripod}
